\documentclass[12pt]{article}
\usepackage[margin=1in, top=0.5in]{geometry} % adjust top margin
\usepackage{amsmath, parskip}

\begin{document}

\begin{center}
\begin{Large}
\textbf{Chapter 1 Checkpoint}
\end{Large}
\end{center}

\vspace{1cm}

(1) Given an account governed by compound interest, what is the relationship between the following parameter? In other words, if I gave you any one of these could you find any other one? \[i, d, v, i^{(m)}, d^{(m)}, \delta\]


\vspace{1cm}


(2) Where does the following expressions move $\$5,000$ (example could be: forward 3 years or backwards 2 months etc.)

\begin{enumerate}
    \item[(a)] $5000 (1+i)^{-4}$
    \item[(b)] $\frac{5000}{(1-d)^3}$
    \item[(c)] $5000(1-\frac{d^{(12)}}{12})^4$
    \item[(d)] $5000(1-\frac{i^{(2)}}{2})^{-8}$
    \item[(e)] $5000v^{-6}$ 
\end{enumerate}

\vspace{1cm}

(3) Would you rather (from an investment percentage) an account governed by...

\begin{enumerate}
    \item[(a)] $i=7\%$ or $v=93\%$
    \item[(b)] $i^{(12)}=1\%$ or $d=12\%$
    \item[(c)] $i=5\%$ or $\delta=5\%$
    \item[(d)] $\frac{i^{(4)}}{4}=8\%$ or $d^{(4)}=2\%$
    \item[(e)] $v=90\%$ or $i^{(2)}=5\%$
\end{enumerate}


\vspace{1cm}

(4) Given an account governed by simple interest of rate $r$, find an expression for the following

\begin{enumerate}
    \item[(a)] The discount rate from year $3$ to year $6$
    \item[(b)] The interest rate from year $2$ to year $4$ 
    \item[(c)] The interest rate from year $3$ to year $5$
    \item[(d)] The force of interest at the start of the 3rd year
    \item[(e)] The amount of interest earned from year 2 to year 7
\end{enumerate}




\end{document}


